\newcommand{\tipoRequerimento}{Contestação do marco temporal de devolução e revisão do período faturado}

\section{DOS FATOS E DOS DIREITOS}

\begin{enumerate}[resume]
\item A \nomeConcessionaria{} julgou procedente a contestação apresentada quanto à Unidade Consumidora \unidadeConsumidora{}, reconhecendo o descabimento da cobrança do custo de disponibilidade para a classe de Iluminação Pública. Contudo, fundamentou indevidamente a limitação temporal do período de devolução ao marco de 31/03/2022, sob o argumento de que tal regra teria passado a vigorar apenas com a entrada em eficácia do art. 668 da REN nº 1.000/2021.
\item Ocorre que a vedação à cobrança de custo de disponibilidade para a classe de iluminação pública não é inovação da REN nº 1.000/2021, já sendo expressamente proibida pela regulação anterior, nos termos do art. 24, § 2º, da Resolução Normativa ANEEL nº 414/2010 ("Não se aplica a cobrança pelo custo de disponibilidade definida no art. 98 no faturamento individual de um ponto de iluminação pública"), dispositivo este sucedido e mantido pelo art. 290, § 1º, inciso I, da REN nº 1.000/2021.
\item Conforme disposto no item 35 da Nota Técnica nº 136/2024/ANEEL, nos casos de condutas continuadas da distribuidora que resultem em faturamento incorreto e que abranjam períodos regidos tanto pela Resolução Normativa nº 414/2010 quanto pela nº 1.000/2021, a análise da irregularidade e o consequente dever de devolução devem ser efetuados de forma unificada, sem interrupção ou limitação arbitrária à data de mudança regulatória.
\item No caso em análise, resta evidente que a cobrança indevida perdura continuadamente desde a vigência da REN nº 414/2010. A recusa da concessionária em restituir os valores cobrados em período anterior a 31/03/2022 descumpre o art. 323 da REN nº 1.000/2021, o Despacho ANEEL nº 2.006/2024 e o entendimento pacificado da Agência Reguladora, caracterizando retenção indevida de valores pagos a maior por erro exclusivo da distribuidora.
\end{enumerate}

\section{DOS PEDIDOS}

\begin{enumerate}[resume]
\item Ante o exposto, requer-se:
    \begin{enumerate}
    \item O afastamento sumário do marco temporal de 31/03/2022 estabelecido pela distribuidora, reconhecendo-se que a proibição de cobrança do custo de disponibilidade para iluminação pública vigora desde a REN nº 414/2010 (art. 24, § 2º) e prossegue na REN nº 1.000/2021 (art. 290, § 1º, I), aplicando-se o entendimento da Nota Técnica nº 136/2024/ANEEL;
    \item A revisão do memorial de cálculo para abranger todo o período de faturamento incorreto desde a vigência da REN nº 414/2010, respeitados os prazos prescricionais regulamentares e a devolução dos valores pagos a maior acrescidos de atualização monetária, juros e devolução em dobro, nos termos do art. 323 da REN nº 1.000/2021 e do Despacho ANEEL nº 2.006/2024.
    \end{enumerate}
\end{enumerate}